\section{matrix element of $e^\frac{(x-X)^2}{2\sigma_\alpha^2}$ in the displaced harmonic oscillator basis}
\label{sec:osc-Vnm-displaced-basis}
\subsection{definition}
consider the displaced  harmonic-oscillator eigenfunctions
\begin{align}
    u_n(x)
    =
    \sqrt\frac{\alpha}{\sqrt{\pi}2^n n!}
    H_n(\alpha x)
    \exp\left(-\frac{\alpha^2x^2}{2}\right),
\end{align}
where $\alpha = \sqrt{M\omega/\hbar}$.
we want to calculate the matrix elements
\begin{align}
    V_{nm}(X,\sigma)
    =
    \left\langle u_n\left|V_0
    \exp\left[-\frac{(x-X)^2}{2\sigma^2}\right]
    \right|u_m\right\rangle .
\end{align}
explicitly,
\begin{align}
    V_{nm}
    =
    V_0 \frac{\alpha}{\sqrt{\pi}\sqrt{2^{n+m}n!m!}}
    \int_{-\infty}^{\infty} dx\,
    H_n(\alpha x)H_m(\alpha x)
    \exp\left(-\alpha^2 x^2\right)
    \exp\left[-\frac{(x-X)^2}{2\sigma^2}\right].
\end{align}
it is convenient to introduce
\begin{align}
    D=\frac{1}{\alpha^2}+2\sigma^2.
\end{align}

\subsection{generating-function representation}
the generating function for the Hermite polynomials is
\begin{align}
    \exp(2zt-t^2)
    =
    \sum_{n=0}^{\infty}H_n(z)\frac{t^n}{n!},
\end{align}
or equivalently
\begin{align}
    H_n(z)
    =
    \left.
    \frac{\partial^n}{\partial t^n}
    \exp(2zt-t^2)
    \right|_{t=0}.
\end{align}
using two auxiliary variables $s$ and $t$,
\begin{align}
    H_n(\alpha x)
     & =
    \left.
    \partial_s^n
    \exp\left({2\alpha sx}-s^2\right)
    \right|_{s=0}, \\
    H_m(\alpha x)
     & =
    \left.
    \partial_t^m
    \exp\left(2\alpha tx-t^2\right)
    \right|_{t=0}.
\end{align}
the Gaussian integral that appears after inserting these representations is
\begin{align}
    \frac{\alpha}{\sqrt{\pi}}
    \int_{-\infty}^{\infty}dx\,
    \exp\left[
    -{\alpha^2 x^2}
    -\frac{(x-X)^2}{2\sigma^2}
    +2\alpha x(s+t)
    -s^2-t^2
    \right].
\end{align}
performing the integral gives the generating function
\begin{align}
    \boxed{
        \mathcal G(s,t)
        =
        \sqrt{\frac{2\sigma^2}{D}}\,
        \exp\left(-\frac{X^2}{D}\right)
        \exp\left[
            \frac{2X}{\alpha D}(s+t)
            -\frac{1}{\alpha^2D}(s^2+t^2)
            +\frac{4\sigma^2}{D}st
            \right].
    }
\end{align}
therefore,
\begin{align}
    \boxed{
        V_{nm}
        =
        V_0  \frac{1}{\sqrt{2^{n+m}n!m!}}
        \left.
        \frac{\partial^{n+m}}
        {\partial s^n\partial t^m}
        \mathcal G(s,t)
        \right|_{s=t=0}.
    }
\end{align}

\subsection{explicit finite-sum expression}
define
\begin{align}
    p=\frac{2X}{\alpha D},
    \qquad
    q=-\frac{1}{\alpha^2 D},
    \qquad
    r=\frac{4\sigma^2}{D}.
\end{align}
then
\begin{align}
    \mathcal G(s,t)
    =
    C\exp\left[p(s+t)+q(s^2+t^2)+rst\right],
\end{align}
where
\begin{align}
    C=
    \sqrt{\frac{2\sigma^2}{D}}
    \exp\left(-\frac{X^2}{D}\right).
\end{align}
expanding all exponentials gives
\begin{align}
    \boxed{
        V_{nm}
        =  V_0      \sqrt{\frac{2\sigma^2}{D}}\,
        e^{-X^2/D}
        \frac{\sqrt{n!m!}}{2^{(n+m)/2}}
    \sum_{\substack{i,j,k\ge0 \\2i+k\le n\\2j+k\le m}}
        \frac{
            q^{i+j}r^k
            p^{n+m-2i-2j-2k}
        }{
            i!\,j!\,k!\,
            (n-2i-k)!\,
            (m-2j-k)!
        }.
    }
\end{align}


\subsection{recurrence relations}
for numerical calculations it is useful to avoid repeated evaluation of
Hermite polynomials or the explicit finite sum. define the unnormalized
coefficients
\begin{align}
    A_{nm}
    =
    \left.
    \partial_s^n\partial_t^m\mathcal G(s,t)
    \right|_{s=t=0}.
\end{align}
they are related to the normalized matrix elements by
\begin{align}
    V_{nm}
    =
    V_0 \frac{A_{nm}}{\sqrt{2^{n+m}n!m!}}.
\end{align}
since
\begin{align}
    \frac{\partial\mathcal G}{\partial s}
    =
    (p+2qs+rt)\mathcal G,
\end{align}
application of $\partial_s^n\partial_t^m$ followed by setting $s=t=0$
gives
\begin{align}
    \boxed{
        A_{n+1,m}
        =
        pA_{nm}
        +
        2qnA_{n-1,m}
        +
        rmA_{n,m-1}.
    }
\end{align}
similarly,
\begin{align}
    \boxed{
        A_{n,m+1}
        =
        pA_{nm}
        +
        2qmA_{n,m-1}
        +
        rnA_{n-1,m}.
    }
\end{align}
using the normalization relation, the first recurrence becomes
\begin{align}
    \boxed{
        \sqrt{n+1}\,V_{n+1,m}
        =
        \frac{p}{\sqrt{2}}V_{nm}
        +
        q\sqrt{n}\,V_{n-1,m}
        +
        \frac{r}{2}\sqrt{m}\,V_{n,m-1}.
    }
\end{align}
in terms of the physical parameters,
\begin{align}
    \boxed{
        \sqrt{n+1}\,V_{n+1,m}
        =
        \frac{\sqrt{2}\,X}{\alpha D}V_{nm}
        -
        \frac{1}{\alpha^2 D}\sqrt{n}\,V_{n-1,m}
        +
        \frac{2\sigma^2}{D}\sqrt{m}\,V_{n,m-1}.
    }
\end{align}
since the operator is real and the oscillator eigenfunctions may be chosen
real,
\begin{align}
    \boxed{
        V_{nm}=V_{mn}.
    }
\end{align}
the corresponding recurrence in the second index is
\begin{align}
    \boxed{
        \sqrt{m+1}\,V_{n,m+1}
        =
        \frac{\sqrt{2}\,X}{\alpha D}V_{nm}
        -
        \frac{1}{\alpha^2 D}\sqrt{m}\,V_{n,m-1}
        +
        \frac{2\sigma^2}{D}\sqrt{n}\,V_{n-1,m}.
    }
\end{align}

\subsubsection{initialization and construction of the matrix}
the starting value is
\begin{align}
    V_{00}
    =
    V_0 \sqrt{\frac{2\sigma^2}{D}}
    \exp\left(-\frac{X^2}{D}\right).
\end{align}
the first element in the first column is
\begin{align}
    V_{10}
    =
    \frac{\sqrt{2}\,X}{\alpha D}V_{00}.
\end{align}
setting $m=0$ in the recurrence gives a simple three-term recurrence for
the first column:
\begin{align}
    \boxed{
        V_{n+1,0}
        =
        \frac{1}{\sqrt{n+1}}
        \left[
            \frac{\sqrt{2}\,X}{\alpha D}V_{n0}
            -
            \frac{1}{\alpha^2 D}\sqrt{n}\,V_{n-1,0}
            \right].
    }
\end{align}
the first row follows from symmetry,
\begin{align}
    V_{0m}=V_{m0}.
\end{align}
once this boundary is known, the complete matrix can be generated using
\begin{align}
    \boxed{
        V_{n+1,m}
        =
        \frac{1}{\sqrt{n+1}}
        \left[
            \frac{\sqrt{2}\,X}{\alpha D}V_{nm}
            -
            \frac{1}{\alpha^2 D}\sqrt{n}\,V_{n-1,m}
            +
            \frac{2\sigma^2}{D}\sqrt{m}\,V_{n,m-1}
            \right].
    }
\end{align}
for $X=0$, the recurrence preserves the parity selection rule
\begin{align}
    V_{nm}=0
    \qquad\text{for }n+m\text{ odd}.
\end{align}
\subsection{simple checks}
for the ground state,
\begin{align}
    \boxed{
        V_{00}
        =
        V_0 \sqrt{\frac{2\sigma^2\alpha^2}{1+2\sigma^2\alpha^2}}\,
        \exp\left[
            -\frac{\alpha^2X^2}{1+2\sigma^2\alpha^2}
            \right].
    }
\end{align}
the first off-diagonal element is
\begin{align}
    \boxed{
        V_{10}=V_{01}
        =
        \frac{\sqrt{2}\,X}{\alpha D}V_{00}.
    }
\end{align}
the element $V_{11}$ is
\begin{align}
    \boxed{
        V_{11}
        =
        \left[
            \frac{2X^2}{D^2\alpha^2}
            +
            \frac{2\sigma^2}{D}
            \right]V_{00}.
    }
\end{align}
for a centered Gaussian, $X=0$, parity implies
\begin{align}
    \boxed{
        V_{nm}=0
        \qquad\text{if }n+m\text{ is odd}.
    }
\end{align}
in the broad-Gaussian limit,
\begin{align}
    \lim_{\sigma\to\infty}
    \exp\left[-\frac{(x-X)^2}{2\sigma^2}\right]=1,
\end{align}
and consequently
\begin{align}
    \boxed{
        \lim_{\sigma\to\infty}V_{nm}=V_0\delta_{nm}.
    }
\end{align}


